\section{The three coefficient lattices and the normalized neighbor map}
\label{sec:edge-lattices}

\status{The integral coefficient-lattice comparisons retained in [JR12] and [LW4] are proved here for the source's polynomial model. This includes full rank, change of lattice coordinates, and the neighbor-map inclusion. An elementary polynomial argument replaces the source's semistandard-basis route to the weight estimate. The construction of the quantitative local tests remains separate.}

\subsection{An integral polynomial model and its scaling}

Let $W$ be a DVR with fraction field $Q$, uniformizer $p$, and residue field of size $q=p^{n'}$. Let $\mathcal O$ be a DVR with fraction field $K$, containing the images of $W$ under embeddings $\tau_s=\iota\sigma^s$, $0\le s<n'$. Assume $w\in\mathcal O$ satisfies $w^2=p\ne0$. Put $k=m+1$. At each $s$ use the decreasing integer weight
\[
 \lambda_s=
 \begin{cases}
   \lambda^0_s,&s\text{ even},\\
   (-\lambda^0_{s,k},\ldots,-\lambda^0_{s,1}),&s\text{ odd},
 \end{cases}
 \qquad
 A=\sum_s\sum_{i=1}^{r}\lambda_{s,i}.
\]
Here $0\le r\le k$, and $A$ is an integer of either sign. These are exactly the source's \texttt{lamAt} and \texttt{Ar}; the reversal at odd $s$ is part of the definition. \textsf{[EL1]}

For a decreasing weight $\lambda$, write $e=\lambda_k$ and $\nu_i=\lambda_i-e\ge0$. Let $\mathcal D_\nu(R)$ consist of polynomials in a $k\times k$ matrix $X$, with column multidegree $\nu$, invariant under every substitution
$C_j\mapsto C_j+uC_i$ with $i<j$. Invariance means a polynomial identity in $u$, over the coefficient ring $R$. Put $T(g)f(X)=f(gX)$, so $T(g)T(h)=T(hg)$, and set
\[
 \rho_\lambda(g)=\det(g)^{-e}T(g^{-1}),\qquad
 H_\lambda(g\mathcal O^k)=\rho_\lambda(g)\mathcal D_\nu(\mathcal O)
 \subset \mathcal D_\nu(K).
 \tag{EL.1}
\]
The column conditions are finitely many linear equations on the finite free module of polynomials of column degree $\nu$. Thus $\mathcal D_\nu(\mathcal O)$ is finite free, and clearing denominators gives
$K\mathcal D_\nu(\mathcal O)=\mathcal D_\nu(K)$.
Row substitutions commute with column additions. Hence $\rho_\lambda(U)$ preserves $\mathcal D_\nu(\mathcal O)$ for $U\in\operatorname{GL}_k(\mathcal O)$, including when $e<0$, since $\det U$ is a unit. This proves that (EL.1) is independent of the basis $g$ and is a full lattice. It also proves
$H_\lambda(a\Lambda)=\rho_\lambda(a)H_\lambda(\Lambda)$.
\textsf{[EL2]}

\begin{lemma}[The needed weight bound without tableaux]
\label{lem:edge-weight}
For $f\in\mathcal D_\nu(K)$, every monomial of $f$ has total degree at most
$N_r=\sum_{i=1}^r\nu_i$ in the first $r$ rows. Consequently, for
$D_w=\operatorname{diag}(w^{-1}I_r,I_{k-r})$,
\[
 w^{\,\sum_{i=1}^r\lambda_i}H_\lambda(g\mathcal O^k)
 \subseteq H_\lambda(gD_w\mathcal O^k).
 \tag{EL.2}
\]
\end{lemma}
\begin{proof}
The cases $r=0,k$ are immediate. Otherwise write the generic matrix as
$X=\left(\begin{smallmatrix}B&C\\ E&F\end{smallmatrix}\right)$, with $B$ of size $r$.
In the polynomial ring localized at $\det B$, column invariance gives
\[
 f\!\left(\begin{pmatrix}zI_r&0\\0&I\end{pmatrix}X\right)
 =
 f\!\left(\begin{pmatrix}zB&0\\ E&F-EB^{-1}C\end{pmatrix}\right).
 \tag{EL.3}
\]
Indeed, multiply on the right by
$\left(\begin{smallmatrix}I&-B^{-1}C\\0&I\end{smallmatrix}\right)$.
Its coefficients are independent of $z$. On the right, every entry involving $z$ lies in the first $r$ columns, whose total prescribed degree is $N_r$.
Localization is injective, so the degree bound holds before localization as well. For integral $f$, it follows term by term that
$w^{N_r}T(D_w)f\in\mathcal D_\nu(\mathcal O)$.
Finally $\det D_w=w^{-r}$ and $\sum_{i\le r}\lambda_i=N_r+re$.
Applying the membership criterion
$h\in H_\lambda(g\mathcal O^k)\iff\det(g)^eT(g)h\in\mathcal D_\nu(\mathcal O)$
therefore gives (EL.2), with no sign restriction on $e$ or on $\sum_{i\le r}\lambda_i$.
\textsf{[EL3--EL4]}
\end{proof}

\subsection{Edges, midpoint lattices, and integral polars}

An edge $e=(L^-,L^+)$ consists of full $W$-lattices with
$pL^+\subseteq L^-\subseteq L^+$ and $|L^+/L^-|=q^r$.
Smith normal form gives, after reordering and absorbing units into basis vectors,
\[
 L^-=gW^k,\qquad
 L^+=g\operatorname{diag}(p^{-1}I_r,I)W^k.
 \tag{EL.4}
\]
To see the count, the elementary divisors divide the uniformizer $p$, so each is a unit or a unit times $p$. The quotient cardinality forces exactly $r$ of the second kind.
Writing $g_s=\tau_s(g)$, the extended source, midpoint, and target lattices are
\[
 g_s\mathcal O^k,\qquad
 \Lambda_s(e):=\tau_s(L^-)\mathcal O+w\,\tau_s(L^+)\mathcal O
      =g_sD_w\mathcal O^k,\qquad
 g_sD_w^2\mathcal O^k .
 \tag{EL.5}
\]
This is the source's \texttt{LeAt}: its factor $w^{-1}p$ equals $w$.
\textsf{[EL5]}

Use separate polynomial variables at different $s$ and the ambient space
$V=\bigotimes_s\mathcal D_{\nu_s}(K)$, where $\nu_{s,i}=\lambda_{s,i}-\lambda_{s,k}$.
This is the source's literal \texttt{Vsp}, with full lattices
\[
 H(L)=\bigotimes_s H_{\lambda_s}(\tau_s(L)\mathcal O),
 \qquad
 H_e=\bigotimes_s H_{\lambda_s}(\Lambda_s(e)).
\]
Tensoring (EL.2) and then applying it a second time gives
\[
 w^A H(L^-)\subseteq H_e,\qquad
 w^A H_e\subseteq H(L^+).
 \tag{EL.6}
\]
For a full lattice $M\subset V$, let
$M^\vee=\{\phi\in V^*:\phi(M)\subseteq\mathcal O\}$; for a lattice in $V^*$ use the analogous polar in $V$.
An $\mathcal O$-basis and its coordinate functionals show
$(M^\vee)^\vee=M$.
Set
\[
 R_e=w^{-A}H_e,\qquad R'_e=w^{-A}H_e^\vee.
\]
Equation (EL.6) gives
$H(L^-)\subseteq R_e$ and $H(L^+)^\vee\subseteq R'_e$.
Moreover the pairing
\[
 \beta_e(x,\phi)=\phi(p^A x),\qquad
 \beta_e(w^{-A}b_i,w^{-A}b_j^*)=\delta_{ij},
 \tag{EL.7}
\]
is perfect over $\mathcal O$. Here the sole scalar calculation is
$w^{-A}p^Aw^{-A}=1$, valid for every integer $A$ because $w^2=p\ne0$.
In particular,
$p^A R_e=(R'_e)^\vee$.
\textsf{[EL6]}

\subsection{The three slots and the neighbor operator}

For an integer $t\ge0$, define the ball slot, the target test lattice, and the dual slot by
\[
 B_t(L)=w^{-t}H(L)\cap\bigcap_{e:\,e^-=L}R_e,\qquad
 T_t(L)=w^{-t}H(L)^\vee\cap\bigcap_{e:\,e^+=L}R'_e,\qquad
 P_t(L)=T_t(L)^\vee .
 \tag{EL.8}
\]
Thus the three coefficient slots are $H(L),B_t(L),P_t(L)$;
$T_t(L)$ itself is a lattice in the dual vector space.
The preceding inclusions imply
\[
 \boxed{\quad
 w^tH(L)\subseteq P_t(L)\subseteq H(L)
       \subseteq B_t(L)\subseteq w^{-t}H(L).
 \quad}
 \tag{EL.9}
\]
Indeed $H(L)^\vee\subseteq T_t(L)\subseteq w^{-t}H(L)^\vee$, and taking polars gives the first two inclusions; the other two follow directly from (EL.8).
Every slot is full and finite free, of the same rank $\dim_K V$: each is sandwiched between two full lattices over the DVR.
This proves all four comparisons in \texttt{RestLocChain}, including the two multiplication-by-$w^t$ statements. \textsf{[EL7]}

For any edge $e=(L^-,L^+)$, (EL.7)--(EL.8) give
\[
 \boxed{\qquad p^A B_t(L^-)\subseteq P_t(L^+).\qquad}
 \tag{EL.10}
\]
The proof is only evaluation: $x\in B_t(L^-)$ lies in $R_e$, while every $\phi\in T_t(L^+)$ lies in $R'_e$, so $\phi(p^Ax)\in\mathcal O$.
The same argument gives
$p^A(w^cR_e)\subseteq w^cP_t(L^+)$ for $c\ge0$.
It therefore also transports the prescribed errors at unwanted neighboring targets, as required in the joint-rank cancellation. \textsf{[EL10]}

For clarity, these constructions respect the group action used on sections.
The identity after (EL.1), applied to all $\tau_s(a)$, sends $H(L)$ and $H_e$ to $H(aL)$ and $H_{ae}$.
On dual vectors use the contragredient action $\phi\mapsto\phi\circ\rho(a)^{-1}$.
It follows by taking intersections and polars that all three slots transform by $\rho(a)$.
Scalar character values in $\mathcal O^\times$ preserve them as well. \textsf{[EL8]}

Finally, on a tuple of lattices $\mathbf L$, the source's neighbor operator is precisely
\[
 (Y_vF)(\mathbf L)
  =p^{A_v}\sum_{L'\,\to\, L_v}
       F(L_1,\ldots,L_{v-1},L',L_{v+1},\ldots).
 \tag{EL.11}
\]
The sum is finite: its possible $L'$ lie between $pL_v$ and $L_v$ and hence correspond to subspaces of the finite vector space $L_v/pL_v$.
Tensor (EL.10) with the unchanged slots at the other places and add the terms. This proves
$Y_v(\text{ball at }v)\subseteq(\text{dual at }v)$.
For equivariant sections, multiplication by $a$ bijects the incoming edges, so reindexing (EL.11) proves equivariance of the result. Sections extended by zero off lattice tuples remain zero there: either there is no edge at $v$, or the unchanged bad coordinate makes every summand vanish.
These observations prove \texttt{RestNat} and its assembly with the coefficient comparisons. They require neither positivity of $A_v$ nor integrality of $p^{A_v}$ by itself. \textsf{[EL9]}

The integral polynomial component of [JR12]/[LW4] is thus discharged. What remains elsewhere is the quantitative production and separation of the local test vectors, the finite-orbit and stabilizer input for the global section space, and the other arithmetic and analytic steps of the RH/GRH argument.
